\section{Alternating infinite behavior: parity strategies}
\label{sec:parity}

\subsection{The source of the additional power}

Suppose a protocol must keep granting requests while an adversary chooses when
to issue them. A finite-state invariant may establish that no illegal grant
occurs. It does not establish that requests are eventually answered, and a
termination rank is inappropriate because the intended protocol runs forever.
An infinite game packages the missing alternation and recurrence information.
Parity games are a standard vehicle for such reasoning; Oink provides a mature
implementation of several solvers~\cite{oink}. Here we deliberately separate
solving from a particularly simple certificate for a supplied strategy.

\begin{definition}[Finite max-parity arena]
An arena consists of a nonempty finite set $S$, a total directed edge relation
$E\subseteq S\times S$, an owner $o:S\to\{0,1\}$, and priorities
$p:S\to\N$. At a vertex owned by player $i$, that player selects the next edge.
An infinite play is won by the parity of the \emph{largest priority occurring
infinitely often}. A positional strategy chooses a successor using only the
current vertex.
\end{definition}

The qualifier ``infinitely often'' and the choice of \emph{largest}, rather than
smallest, must appear in the source bridge and the serialization version. A
solver using min-parity can be adapted only through an explicit,
parity-preserving conversion. Dead ends are rejected: declaring them losing,
winning, or self-looping would otherwise change the problem.

\subsection{A local certificate with no trusted game solver}

For player $i$, the certificate supplies a region $W_i$ and a strategy
$\sigma_i$ on the vertices of $W_i$ that it owns. The checker first establishes
\begin{align}
  o(v)=i,\ v\in W_i &\Longrightarrow
     (v,\sigma_i(v))\in E\ \land\ \sigma_i(v)\in W_i,\label{eq:own}\\
  o(v)\ne i,\ v\in W_i &\Longrightarrow
     \forall w,\ (v,w)\in E\Rightarrow w\in W_i.\label{eq:opp}
\end{align}
Let $H_i$ retain the selected edge at owned vertices and \emph{all} edges at
opponent vertices. This graph represents every play consistent with the
proposed strategy. Checking only the edges the solver happened to explore would
be unsound.

For every priority $q$ occurring in $W_i$ whose parity is unfavorable to $i$,
form the threshold graph
\[
 H_{i,\le q}=H_i[\{v\in W_i:p(v)\le q\}].
\]
The certificate gives a natural-number rank $r_q$ on that graph. Every retained
edge $v\to w$ with both priorities at most $q$ must satisfy
\begin{equation}
 r_q(v)\ \ge\ r_q(w)+\one{p(v)=q}.
 \label{eq:parity-rank}
\end{equation}
There is no condition across an edge passing through a priority larger than
$q$. The rank may reset there. This is exactly what a single global decreasing
rank cannot express.

The implementation stores one rank array per unfavorable threshold, with
irrelevant entries fixed to zero and every relevant entry at most $|S|$. It
also supplies both regions and checks that they partition the entire arena.
For a goal about one initial state, a single suitable region is mathematically
enough; the stronger partition contract makes the delivered solver's complete
answer easy to test.

\begin{theorem}[Soundness of threshold-local ranks]
\label{thm:parity}
If \eqref{eq:own}, \eqref{eq:opp}, and \eqref{eq:parity-rank} hold, every play
starting in $W_i$ and following $\sigma_i$ is won by player $i$.
\end{theorem}
\begin{proof}
The closure checks keep every vertex of the play in $W_i$. Since $S$ is finite,
at least one priority occurs infinitely often. Let $q$ be the largest such
priority. All priorities greater than $q$ occur only finitely many times, so
some tail of the play lies in $H_{i,\le q}$. If $q$ were unfavorable, its
certificate rank would never increase along this tail and would strictly
decrease on every departure from a $q$-vertex. There are infinitely many such
departures. A natural number admits only finitely many strict decreases, a
contradiction. Thus $q$ has parity $i$.
\end{proof}

\begin{proposition}[Completeness of the certificate for a winning strategy]
\label{prop:rank-complete}
Fix a positional strategy that wins from every vertex in a closed region $W_i$.
For each unfavorable $q$, ranks satisfying \eqref{eq:parity-rank} exist with
values at most $|S|$.
\end{proposition}
\begin{proof}
The threshold graph has no directed cycle containing a $q$-vertex. Otherwise
an opponent could follow that cycle forever; all owned edges on it already
agree with the fixed strategy. The largest recurrent priority would be $q$,
contradicting winningness. Give each threshold edge weight
$\one{p(v)=q}$. Every directed cycle therefore has weight zero. The largest
weight of a finite walk starting at $v$ is finite: removing repeated-vertex
cycles never changes its weight, so a maximizing walk has weight at most
$|S|$. This largest weight is a rank satisfying the required inequalities.
\end{proof}

Together with positional determinacy of finite parity games, the proposition
explains why the certificate language is sufficient for complete finite-arena
solving. It does \emph{not} make the recursive search polynomial, and it is not
a claim to have implemented the lexicographic small-progress-measures algorithm.
The certificate is intentionally a simple post-processing object for a chosen
strategy. Standard solver and progress-measure families are documented in
Oink~\cite{oink}.

\subsection{Concrete search algorithm}

The proposer uses recursive Zielonka solving. In a subarena $U$, let $q$ be the
largest priority, let $i=q\bmod2$, and let $T$ contain the vertices of priority
$q$. Compute the attractor $A=\operatorname{Attr}_i(T,U)$: repeatedly add an
$i$-vertex with a successor already in $A$, or an opponent vertex all of whose
subarena successors lie in $A$. Store the chosen move whenever an owned vertex
is added.

Solve $U\setminus A$ recursively. If the opponent's winning region is empty,
player $i$ wins $U$ using the attractor strategy and the recursive strategy.
Otherwise attract the opponent to its recursively winning region, remove this
second attractor, solve the remainder, and combine the two strategies.
The implementation checks that recursive subarenas remain total; a malformed
recursion is a diagnostic failure, never negative evidence.

\begin{lstlisting}
solve(U):
  if U is empty: return empty winning regions and strategies
  q := maximum priority in U
  i := q mod 2
  A, reachA := attract(i, vertices of priority q, U)
  W, sigma := solve(U minus A)
  if W[1-i] is empty:
    return combine(i wins A, reachA, W, sigma)
  B, reachB := attract(1-i, W[1-i], U)
  V, tau := solve(U minus B)
  return combine(1-i wins B, reachB, sigma[1-i], V, tau)
\end{lstlisting}

After strategies are found, the prototype computes rank arrays by repeated
relaxation of \eqref{eq:parity-rank}, starting at zero. A nonconvergent positive
cycle causes rejection. A production emitter should instead compute strongly
connected components of each threshold graph and longest weighted paths in
its condensation; this changes rank extraction, not checking.

Writing $n=|S|$, $m=|E|$, and $d$ for the number of distinct unfavorable
thresholds, the rank certificate occupies $O(dn\log n)$ bits apart from the
strategy and region encoding. Its mathematical local checker performs $O(dm)$
rank comparisons and $O(n+m)$ closure work. The supplied Python implementation
uses straightforward list/set membership rather than claiming an optimized
realization of every asymptotic bound. Its rank emitter uses at most $O(dnm)$
relaxation work. Recursive Zielonka search retains its exponential worst-case
behavior; an Oink adapter would be the appropriate scaling experiment.

\subsection{Two semantic regression examples}

Consider the forced cycle
\[
 a\ (p=1)\ \longrightarrow\ b\ (p=2)\ \longrightarrow\ a.
\]
Player 0 wins because priority 2 recurs infinitely often. There is no natural
rank strictly decreasing around the whole cycle. For the sole unfavorable
threshold $q=1$, the threshold graph has the vertex $a$ but no edge, so the
zero rank suffices. This is not evasion: visiting priority 2 is precisely what
makes unbounded repetition acceptable.

By contrast,
\[
 a\ (p=2)\ \longrightarrow\ b\ (p=1),\qquad b\longrightarrow b
\]
is won by player 1. The largest priority seen anywhere is 2, but the largest
one seen infinitely often is 1. Both examples are executed regressions. Another
regression constructs a strategy certificate for an arena with one adversarial
edge removed and verifies that it is rejected against the original arena.

\subsection{How the strategy becomes a Lean witness}

The source result should have the form ``there exists a legal controller such
that every legal environment run satisfies the parity condition.'' A table
$\sigma_i$ is an executable witness once a proved finite enumeration maps source
states and actions to the arena. The proof uses Theorem~\ref{thm:parity}; it
must not identify an arbitrary source function with a vertex table by pretty
printing or by sampling its behavior.

A finite arena may be an abstraction of an infinite-state source program.
Then a winning abstract controller transfers only with a proved alternating
simulation in the correct direction: the controller's abstract choices must be
realizable, and every concrete adversarial move must be covered. An
opponent-winning abstract strategy need not refute concrete controller
existence if the abstraction introduced spurious adversarial moves. Negative
transfer consequently needs an exact model or a separately proved reverse
bridge. This distinction is not encoded by a bare Boolean ``abstract model
verified'' flag.

Spot can propose temporal automata and manipulate parity games~\cite{spot}.
However, checking a strategy in an automaton does not check that the automaton
recognizes the original temporal formula. The first release should accept
explicit, already interpreted parity arenas. A verified temporal translation
can be added later without putting the translator into the trusted boundary.
